\documentclass[10pt,a4paper]{amsart}
\usepackage[margin=19mm]{geometry}
\usepackage{amsmath,amssymb,mathtools}
\usepackage[scr=rsfs,cal=euler]{mathalfa}
\usepackage{xfrac}
\usepackage[unicode,hidelinks,bookmarks=false]{hyperref}
\newcommand{\ie}{i.e.\ }
\newcommand{\cf}{cf.\ }
\newcommand{\resp}{respectively\ }
\newcommand{\Subsection}{\subsection}
\providecommand{\nobreakdash}{\nobreak\hbox{-}}
\setlength{\emergencystretch}{2em}
\multlinegap0pt
\usepackage[shortlabels]{enumitem}
\usepackage{amssymb}
\newtheorem{lemma}{Lemma}
\newcommand{\e}{\mathrm{e}}
\newcommand{\probP}{\mathbb{P}}
\newcommand{\reals}{\mathbb{R}}
\title{E-variables and tests of randomness for distribution classes}
\author{Georgii Potapov, Yuri Kalnishkan}
\date{Source: arXiv:2603.02492v1 (CC BY 4.0)}
\begin{document}
\maketitle

\noindent\emph{Source and licence.} E-variables and tests of randomness for distribution classes. Version arXiv:2603.02492v1 (CC BY 4.0), distributed by arXiv under the Creative Commons Attribution 4.0 International licence, http://creativecommons.org/licenses/by/4.0/, as recorded in the arXiv licence metadata for this exact version and shown on the abstract page https://arxiv.org/abs/2603.02492v1. Full licence notice: \texttt{licenses/arxiv-2603.02492-CC-BY-4.0.txt}.\par\smallskip

\noindent\textit{Editorial scope.} Counted unit: the article's lemma lem:4-cond-techn-lemma (source0.tex lines 326--350). The article's complete original proof of it is delivered with it (source0.tex lines 656--730, one \texttt{proof} environment of 75 lines, which the article places away from the statement). No mathematics is written, repaired or re-derived: the statement and the proof are the article's own lines, in the article's own notation, and no retained line is substituted or re-flowed.  No further source-local context is needed: every cross-reference of every retained line resolves inside the two delivered spans.  Nothing was omitted from any retained span, so the package carries no omission record.  No retained line contains a citation, so the wrapper carries no bibliography and no bibliography entry.  No retained line contains a figure environment, an image-inclusion command or drawing code, so no image is delivered and the package declares no asset dependency.  Editorial formatting only, in addition: an AMS \texttt{amsart} preamble that restates the article's own declarations and macros used by the retained lines, namely the environments \texttt{lemma} on separate counters, together with the standard packages those lines need.  The retained environments print in this document as Lemma 1 in order of appearance, whereas the article prints the same items with the numbering of its own full document.  The complete native source of the article as distributed in the arXiv e-print for this exact version is bundled unchanged as \texttt{sources/195623/source0.tex}.
\begin{lemma}\label{lem:4-cond-techn-lemma}
    Let $\mathcal{H} = \{\probP_\theta\}_{\theta\in\Theta}$ be a family of distributions on $\reals^n$ with densities $\{p_\theta\}_{\theta\in\Theta}$, parametrised by the elements of $\Theta\subseteq\reals$, and let $\mathcal{S} \subseteq \Theta$ be an at most countable set such that $\operatorname{succ}_{\mathcal{S}}(\theta)=\min\{s\in\mathcal{S}\mid\theta<s\}$ and $\operatorname{pred}_{\mathcal{S}}(\theta)=\max\{s\in\mathcal{S}\mid s<\theta\}$ are well-defined for all $\theta\in\Theta$ expect for possibly those that are greater than maximal or less than minimal element of $\mathcal{S}$. Let $d(\cdot\|\cdot):\Theta\times\Theta\to [0, +\infty]$ be a well-defined function. Let $\hat s:\reals^n\to \mathcal{S}$ be a statistic and $g:\reals^n\to\Theta$ be a function.
    Consider now the following properties:
    \begin{enumerate}[label=(p \roman*)]
        \item\label{prop:loglike-to-d} For every $\theta \in\Theta$, $s \in \mathcal{S}$ and $x\in \hat{s}^{-1}(s) \cap \{y\mid p_\theta(y)>0\}$ we have that
        \[
            \log\frac{p_\theta(x)}{p_{s}(x)} = d(g(x)\|s)-d(g(x)\|\theta);
        \]
        \item\label{prop:near-bound} there is a constant $c' > 0$ such that for any $s \in \mathcal{S}$ and any $x$ for which $x \in \hat s^{-1}(s)$ we have
        \[
            d(g(x)\|s) < c'.
        \]
        \item\label{prop:g-near-s} For all $s\in\mathcal{S}$ and all $x \in \hat s^{-1}(s)$
        \[
            \operatorname{pred}_{\mathcal{S}}(s)\le \operatorname{pred}_{\mathcal{S}}(g(x))\le \operatorname{succ}_{\mathcal{S}}(g(x))\le \operatorname{succ}_{\mathcal{S}}(s).
        \]
        \item\label{prop:d-grows-fast-enough} There is a constant $\alpha >0$ such that for all $\theta_1, \theta_2 \in \Theta, \theta_1 < \theta_)$, it holds that
        \[
            d(\theta_1\|\theta_2) \ge (1+\alpha)\log (k-1), ~~~ d(\theta_2\|\theta_1) \ge (1+\alpha)\log (k-1),
        \]
        for $k = \left| \mathcal{S} \cap (\theta_1, \theta_2) \right|$\footnote{for this statement, we treat the logarithm of non-positive numbers as $-\infty$}.
    \end{enumerate}

    If all the properties \ref{prop:loglike-to-d}--\ref{prop:d-grows-fast-enough} of the data $(\mathcal{H}, \mathcal{S}, \Theta, \hat{s}, g, d)$ hold, then $\mathcal{H}$ is $e$-variable approximable from the net $\mathcal{S}$ and the estimator $\hat s$, and the constant factor $C$ can be taken as $C =\exp(c')\left(7 +  2/\alpha\right)$.
\end{lemma}

\begin{proof}[Proof of lemma \ref{lem:4-cond-techn-lemma}]
    We want to show that $\frac{1}{C}e_{\hat s (X)}(X)$ is an $e$-variable. That is, we need to show that, for all $\theta \in \Theta$
    \[
        \int_{\reals^n} e_{\hat s(x)}(x) \probP_\theta({\rm d}x)\leq C.
    \]
    First, for each $s \in \mathcal{S}$ we denote $\mathcal{X}_s = \hat s^{-1}(s)\subseteq \reals^n$. Now,
    \begin{align*}
        \int_{\reals^n} e_{\hat s(x)}(x) \probP_\theta({\rm d}x) 
        &=
        \sum_{s\in\mathcal{S}} \int_{\mathcal{X}_s}e_s(x) p_{\theta}(x) {\rm d}x
        \\ &=
        \sum_{s\in\mathcal{S}} \int_{\mathcal{X}_s}
            \frac{p_\theta(x)}{p_s(x)}
            e_s(x)p_s(x){\rm d}x
        \\ &\overset{\text{by \ref{prop:loglike-to-d}}}{=}
        \sum_{s\in\mathcal{S}} \int_{\mathcal{X}_s}
            \exp\Big[d(g(x)\|s) - d(g(x)\|\theta)\Big]
            e_s(x)p_s(x){\rm d}x 
        \intertext{denote 
            $C_{s}=C_{s,\theta}   =\sup_{x\in\mathcal{X}_s}d(g(x)\|s) - d(g(x)\|\theta)$}
        &\le\sum_{s\in\mathcal{S}}
            \exp(C_s)\int_{\mathcal{X}_s}
            e_s(x)p_s(x){\rm d}x 
        \\ &\le\sum_{s\in\mathcal{S}}
            \exp(C_s)\underbrace{\int_{\reals^n}e_s(x)p_s(x){\rm d}x}_{\le 1,\textrm{ by def. of $e$-var.}}
        \\ &\le\sum_{s\in\mathcal{S}}
            \exp(C_s).
    \end{align*}
    We now have to estimate the sum $\sum_{s\in\mathcal{S}}
    \exp(C_s).$

    By \ref{prop:near-bound} we have 
    \[
        C_{s,\theta} \leq c' - \inf_{x\in\mathcal{X}_s} d(g(x)\|\theta),
    \]
    and, from non-negativity of $d$, it follows that
    $
        C_s \leq c'.
    $

    Whenever $s \ge \operatorname{succ}_{\mathcal{S}}^{(k+1)}(\theta)$ for some $k$, then $\operatorname{pred}_{\mathcal{S}}(s) \ge \operatorname{succ}_{\mathcal{S}}^{(k)}(\theta)$, and, by \ref{prop:g-near-s} for all $x \in \mathcal{X}_s$:
    \[
        g(x) > \operatorname{pred}_{\mathcal{S}}(g(x)) \ge \operatorname{pred}_{\mathcal{S}}(s) \ge \operatorname{succ}_{\mathcal{S}}^{(k)}(\theta) > \theta,
    \]
    and $\left|\mathcal{S} \cap (g(x), \theta) \right| \ge k$, and, by \ref{prop:d-grows-fast-enough} we have: 
    \[
        d(g(x)\| \theta) \ge (1+\alpha)\log (k-1), ~~\forall x\in\mathcal{X}_s,
    \]
    and so, in this case $C_{s, \theta} \le c' - (1+\alpha)\log (k-1)$.

    Similarly, we get the same bound for $s \le \operatorname{pred}_{\mathcal{S}}^{(k+1)}(\theta)$.
    
    Putting this all together, we get
    \begin{align*}
        \sum_{s\in\mathcal{S}}
            \exp(C_{s,\theta}) &\le
            \left(\sum_{\substack{s \le \operatorname{pred}_{\mathcal{S}}^{(3)}(\theta) \\ s\in\mathcal{S}}}
            +
            \sum_{\substack{\operatorname{succ}_{\mathcal{S}}^{(3)}(\theta) \le s \\ s\in\mathcal{S}}}
            +
            \sum_{\substack{s\in \mathcal{S} \\ \operatorname{pred}_{\mathcal{S}}^{(3)}(\theta) < s \\ s < \operatorname{succ}_{\mathcal{S}}^{(3)}(\theta)}}
            \right) \exp(C_{s,\theta})
            \\
            &\le \sum_{k=1}^\infty \e^{c'-(1+\alpha)\log k} + \sum_{k=1}^\infty e^{c'-(1+\alpha)\log k} + 5\e^{c'} 
            \\
            &= \e^{c'} \left(5 + 2\sum_{k=1}^{\infty} \frac{1}{k^{1+\alpha}}\right) 
            \\
            &\le
            e^{c'} \left(5 + 2\left(1 + \frac{1}{\alpha}\right)\right)
            \\
            &=
            \exp(c')\left(7 + \frac{2}{\alpha}\right),
    \end{align*}
    which concludes the proof.
\end{proof}

\end{document}
